\section{Schéma électrique}
Insérer le schéma de mathis quand il sera fini

\section{Choix de câblage et choix technique}
Cette section détaille le choix de câblage des différents composants de notre système ainsi qque les raisons techniques de ces choix.

\subsection{Ecran LCD} 
L'écran LCD communique avec la Raspberry Pi par le protocole I2C. Nous avons donc utilisé :
\begin{itemize}
    \item le GPIO 2 pour la ligne SDA (Data)
    \item le GPIO 3 pour la ligne SCL (Clock)
\end{itemize}
L'écran est alimenté en 5V depuis la Raspberry Pi.

\subsection{Capteur RGB}  
Le capteur RGB fonctionne également avec le protocole de communication I2C. La Raspberry Pi ne dispose que d'un seul couple de GPIO pour SDA et SCL, nous avons du utilisé des dominos afin de brancher à la fois l'écran LCD et le capteur RGB sur les mêmes lignes I2C.

\subsection{Buzzer}
Le buzzer ne nécessite pas de protocole particulier. Nous l'avons donc connecté sur le GPIO 25.

\subsection{Capteurs de ligne}
Les capteurs de lignes ont été mis sur des GPIO n'ayant pas de particularité c'est à dire le GPIO.... Ce choix permet une lecture simple des signaux numériques.

\subsection{Moteur}  
Les connexions pour le moteur ont été choisies de la manière suivante :
\begin{itemize}
    \item EN1 et EN2 sur les GPIO 13 et 12 pour le contrôle en PWM. Ces 2 GPIO vont permeterre de controler la vitesse des moteurs.
    \item IN1,IN2,IN3,IN4 ont été mit sur des GPIO standards. Ils nous permettront de définir le sens de rotation des moteurs.
\end{itemize}
Cette configuration permet un contrôle indépendant des moteurs et une modulation de la vitesse essentielle pour la navigation précise du robot.

\section{Bilan de puissance et estimation de l'autonomie}
????

\section{Rôles des composants}

\subsection{Raspberry Pi}
La Raspberry Pi est comme le cerveau du robot, elle éxécute les programmes pour contrôler les moteur et elle interprète les données des capteurs.

\subsection{batterie USB-C}
Cette batterie permet d'alimenter la Raspberry Pi. Cette batterie permet d'assurer que la Raspberry Pi reste toujours allumée même si les moteurs sont éteints.

\subsection{Pile}
Les piles vont fournir un courant plus élevé aux moteurs, ces piles sont nécessaires pour transmettre un courant assez élevé aux moteurs sans surcharger la Raspberry Pi.

\subsection{Ecran LCD} 
Le rôle de l'écran LCD est d'assurer l'interface homme-machine du robot. Il affiche en temps réel les informations suivante : les ordres éxécutés (AV, TG, TD), la couleur détectée par le capteur RGB ainsi que l'état du robot (initialisation, arret).
L'écran LCD est un afficheur utilisé par le protocole de communication I2C et à partir d'un module d'extension ( PCF8574) qui va nous faciliter les branchements.
Le protocole I2C permet de réduire le nombre de fils nécessaires en utilisants seulement : SDA (données) et SCL (clock).

\subsection{Capteur RGB}
Le capteur RGB permet au robot de détecter la couleur des pastilles présentes sur le circuit. Au départ il evra détecter une pastille rouge puis il devra détecter les pastilles vertes qui seront les points de passage obligatoire, et pour finir il devra revenir à la pastille rouge. Lorqu'une de ces couleurs est détectée, le robor doit afficher cette information sur l'écran LCD et activer le buzzer.
Ce capteur mesue donc l'intensité de la lumière réfléchie par trois composantes qui sont les trois couleurs primaires, le rouge , le vert et le bleu. Une LED est intégrée à ce module et permet d'éclairer la surface à analyser. Des photodiodes internes mesurent la lumière réfléchie et le capteur convertit ces mesures en valeurs numériques. Ces valeurs sont transmises vers la Rasbpberry Pi via le protocole bus I2C. 
Le programme devra donc comparer les valeurs R, G,B afin de déterminer la couleur dominante détectée.

\subsection{Buzzer}
Le buzzer va permettre d'mettre un bip lorsque le robot atteint un point de passage obligatoire, c'est à dire lorsque le capteur RGB a détecté une pastille rouge ou verte.
Le buzzer est commandée directement par une sortie GPIO de la Raspberry Pi. Lorsque la broche GPIO est mise à l'état haut le buzzer est alimenté et produira en conséquence un son. 

\subsection{Capteurs de ligne }
Les capteurs de ligne permettent au robot de suivre les lignes noires tarcées sur le sol et de détecter les croisements. Ils sont indispensables pour assurer le guidage précis du robot dans la ville.
Chaque capetur fonctionne à l'aide d'une LED infrarouge et d'un phototransistor. Une surface claire réfléchit fortement la lumière alors qu'une ligne sombre va absorbé cette lumière. Selon cette quantité de lumière reçue, le capteur fournira un signal logique (0 ou 1) à la Raspberry Pi. 

\subsection{Moteurs}
Les moteurs assurent les déplacements du robot : avancer en ligne droite,tourner à gauche, tourner à droite et s'arréter.
Le robot utilise deux moteurs à courant continu, un pour chaque roue. Nous allons pouvoir contrôler le sens et la vitesse de chaque moteur via PWM.

\subsection{Interrupteur}
L'interrupteur va nous permettre de couper ou de remettre l'alimentation électrique des moteurs sans affecter la Raspberry Pi. Ce composant va permettre de mettre le robot en sécurité avant l'intervention, de démarrer ou d'arrêter le déplacement du robot facilement et de préserver la Raspberry Pi lorsqu'on manipule le robot.



\section{Organigramme du rôle et du principe de fonctionnement des composants}
?????

\section{Programmation}
Détaillée les codes utilisés

\section{Dificultés rencontrées et pistes d'améliorations}

